\subsubsection{REDES \en{PEER-TO-PEER}}

Redes \en{peer-to-peer} são redes onde cada \en{node} atua como cliente e servidor. Essas redes podem ser de arquitetura estruturada e não estruturada. Nas estruturadas a rede é construída usando um procedimento determinista; um exemplo seria manter uma tabela com identificadores para todos os \en{nodes}, este mesmo identificador atua para categorizar os dados, assim quando um item é acessado na rede, é retornado o identificador do \en{node} responsável para ele. Em uma rede não estruturada os \en{nodes} conhecem uns aos outros de maneira quase aleatória, isso é, cada um mantém uma lista de vizinhos construída de maneira não determinista; a transmissão de dados ocorre da mesma forma \cite{tanenbaum2007dist}.
